\section{Phạm vi đề tài}

\textit{Về mục tiêu chung}, đề tài hướng tới việc xây dựng một giải pháp trọn gói cho hệ thống Thương mại điện tử hiện đại, tích hợp framework xử lý dữ liệu thông minh Florus \cite{vo2024florus}. Mục tiêu cốt lõi là giải quyết bài toán Real-time Personalization và xử lý vấn đề Cold-start thông qua kiến trúc Microservices có tính scalability linh hoạt.

\textit{Về mục tiêu của giai đoạn}, tập trung vào phân tích nghiệp vụ, thiết kế kiến trúc hệ thống, cơ sở dữ liệu và xây dựng khung hạ tầng trên Kubernetes. Cụ thể, giai đoạn này ưu tiên việc thiết lập môi trường vận hành tự động hóa (DevOps/GitOps), đóng gói framework Florus và đảm bảo luồng dữ liệu được lưu thông suốt từ E-commerce App đến các kho dữ liệu.

\textit{Về kết quả mong đợi}, kết thúc giai đoạn Đồ án chuyên ngành, nhóm thực hiện dự kiến đạt được:
\begin{itemize}
    \item \textbf{Về hạ tầng:} Xây dựng thành công Cluster Kubernetes hoàn chỉnh (trên Minikube/Cloud) được quản lý bằng Helm Charts (mô hình Umbrella Chart).
    \item \textbf{Về kiến trúc:} Hiện thực hóa kiến trúc Microservices tách biệt giữa tầng ứng dụng (E-commerce Router) và tầng xử lý dữ liệu (Florus Pipeline).
    \item \textbf{Về dữ liệu:} Thiết lập thành công luồng thu thập dữ liệu (Ingestion) qua Kafka và tổ chức lưu trữ đa tầng (Redis cho Hot data, Neo4j cho Graph data).
    \item \textbf{Về ứng dụng:} Triển khai bản mẫu (Proof of Concept) chứng minh khả năng giao tiếp giữa Router và Florus Pipeline để hiển thị danh sách gợi ý cơ bản.
\end{itemize}

\textit{Về thời gian}, đề tài bao gồm hai giai đoạn chính:
\begin{itemize}
    \item Giai đoạn \textit{Đồ án chuyên ngành} kéo dài trong 14 tuần.
    \item Giai đoạn \textit{Đồ án tốt nghiệp} kéo dài trong 15 tuần.
\end{itemize}

\textit{Về lĩnh vực hướng đến}, đề tài tập trung vào giao thoa giữa \textbf{Kỹ thuật dữ liệu (Data Engineering)} và \textbf{Vận hành hệ thống (DevOps/SRE)}. Cụ thể:
\begin{itemize}
    \item Ứng dụng Container Orchestration (Kubernetes) để quản lý vòng đời ứng dụng Big Data.
    \item Xây dựng hệ thống gợi ý (Recommender Systems) dựa trên phiên (Session-based) và đồ thị tri thức (Knowledge Graph).
    \item Thiết kế hệ thống phân tán chịu tải cao (High Availability Distributed Systems).
\end{itemize}